% Standalone section; requires amsmath, amssymb, amsthm and the usual
% theorem/proposition/lemma environments.  It uses no private macros.
% Additional bibliography key needed: MS (entry below).
% All arithmetic assertions in this section are unconditional at p=37.

\section{The component quotient and its Galois action}
\label{sec:component-quotient-audit}

Put $F_1=\mathbb Q(\mu_p)$, $F_\infty=\mathbb Q(\mu_{p^\infty})$,
and let $\Omega_p$ be the maximal pro-$p$ extension of $F_1$
unramified outside $p$.  Write
\[
 P=\operatorname{Gal}(\Omega_p/F_1),\qquad
 H=\operatorname{Gal}(\Omega_p/F_\infty),\qquad
 \Gamma=\operatorname{Gal}(F_\infty/F_1),\qquad
 \Lambda=\mathbb Z_p[[\Gamma]].
\]
We use closed commutator and normal subgroups throughout.  Choose a
lift of $\Delta=\operatorname{Gal}(F_1/\mathbb Q)$ by the profinite
Schur--Zassenhaus theorem.  Coprime fixed-point lifting gives a
$\Delta$-fixed element of $P$ mapping to a generator of $\Gamma$;
its procyclic closure is a section.  Thus we may use commuting
lifts and write $P=H\rtimes\Gamma$.  Complex conjugation is denoted
by $c$.  For $M=H^{\rm ab}$ write $U=M^-$ and $T=M^+$.

\subsection{An explicit calculation with arbitrary lifts}

The choice of an abelian component does not specify a subgroup to be
killed in a nonabelian group.  The following calculation gives the
exact maps behind this issue.  Let $F$ be free pro-$p$ on $x,y,t$, where
$p$ is odd, and define
\[
 c(x)=x^{-1},\qquad c(y)=y^{-1},\qquad c(t)=t.
\]
Thus $F^{\rm ab}=\mathbb Z_p x\oplus\mathbb Z_p y\oplus\mathbb Z_p t$
has two negative summands and one positive summand.  With
$[x,y]=x^{-1}y^{-1}xy$, both
\[
 a\longmapsto t^a,\qquad
 a\longmapsto(t[x,y])^a
\]
lift the same homomorphism $\mathbb Z_p\to F^{\rm ab}$.
Their normal closures are the respective kernels of the retractions
to the free pro-$p$ group on $x,y$ which send $t$ to $1$ and
$[x,y]^{-1}$.  These are explicit identifications of both quotient
groups, with their different quotient maps.

Let $E_{ij}$ denote a matrix unit and set, in
$\operatorname{UT}_4(\mathbb F_p)$,
\[
 X=I+E_{12},\qquad Y=I+E_{23}+E_{34},\qquad Z=I-E_{13}.
\]
Matrix multiplication gives
\[
 [X,Y]=I+E_{13},\qquad
 [X^{-1},Y^{-1}]=I+E_{13}-E_{14}.
\]
The homomorphism $x\mapsto X$, $y\mapsto Y$, $t\mapsto Z$
therefore kills $t[x,y]$ and sends $c(t[x,y])$ to $I-E_{14}$.
Consequently
\begin{equation}\label{eq:component-raw-lift}
 c\bigl(\langle t[x,y]\rangle_F^{\rm normal}\bigr)
 \ne\langle t[x,y]\rangle_F^{\rm normal}.
\end{equation}
This is a finite quotient certificate for the failure of invariance.

Exact fixed lifts also leave choices.  In the free group on $x,y$ put
\[
 h=[x,y],\qquad d=h^{-1}c(h),\qquad w=h d^{1/2}.
\]
Unit powers are uniquely defined in a pro-$p$ group.  In particular
$c(d)=d^{-1}$, so $c(w)=w$.  The class of $w$ modulo the third lower
central subgroup is the nonzero class of $[x,y]$.  Its nonzero image
is also visible in $\operatorname{UT}_3(\mathbb F_p)$.
Thus $t$ and $tw$ are exact $c$-fixed lifts of the same abelian
generator with distinct $c$-stable normal closures: the quotient
killing $t$ sends $tw$ to $w\ne1$.

\subsection{The actual arithmetic group at $p=37$}

We first give all arithmetic inputs used below.  Iwasawa's
Theorem~2, together with Ferrero--Washington, makes $H$ a free
pro-$37$ group
\cite[Theorem~2, pp.~273--274]{Iwasawa,FW}.
The maximal extension in that theorem agrees with $\Omega_{37}$:
a finite $37$-extension over $F_\infty$ descends to a finite
cyclotomic layer, and its normal closure over $F_1$ remains a
$37$-extension unramified outside $37$.  The group $H$ is
second countable as a closed subgroup of the finitely generated
pro-$37$ group $P$.

The cyclotomic class-group calculation at $37$ gives
\begin{equation}\label{eq:component-p37-class-input}
 X^+=0,\qquad
 X^-=X^{[5]}\simeq\mathbb Z_{37}\quad
 \text{as a $\mathbb Z_{37}$-module}.
\end{equation}
Here the brackets denote the Teichm\"uller character exponent.
One published verification is
\cite[\S1, items~(2)--(3)]{HHO}: the $37$-class group
of $\mathbb Q(\mu_{37^r})$ is cyclic of order $37^r$ for every
$r\ge1$.  Its character is $5=1-32$ modulo $36$.
The direct limit $A_\infty^-$ under extension of ideals is therefore
$(\mathbb Q_{37}/\mathbb Z_{37})^{[5]}$.
For example, after choosing compatible generators for the
surjective norm maps, the extension maps at all sufficiently high
levels are multiplication by $37$ times a unit, since the composite
of norm and extension is multiplication by $37$.
The exact Kummer duality map
\[
 M^+\simeq(A_\infty^-)^\vee(1)
\]
then gives
\begin{equation}\label{eq:component-p37-even}
 T=T_{32}\simeq\mathbb Z_{37},\qquad T_e=0\quad(e\ne32).
\end{equation}
The character transport is explicit: on the dual functional
$f$ the action is $(\delta f)(b)=\omega(\delta)f(\delta^{-1}b)$,
which sends exponent $5$ to $1-5=32$ modulo $36$.
For the Kummer map see
\cite[Theorem~2.3.18]{SharifiCyclotomicNotes}.

The integral odd module is
\begin{equation}\label{eq:component-p37-odd}
 U\simeq\Lambda^{18}.
\end{equation}
We include the two exceptional character calculations in the proof
of this identity.  For an odd exponent $a\ne1,-1$, apply
\cite[Proposition~2.2.6]{WWEOrd} with
$\theta=\omega^{1-a}$.  The proposition identifies $X_\theta=0$
with integral rank-one $\Lambda$-freeness of
$M^{[a],\#}(1)$.  The inversion and twist transport the ring by
$\gamma\mapsto\kappa(\gamma)\gamma^{-1}$, an automorphism; undoing
them preserves freeness and rank.  This proves the assertion in
sixteen of the eighteen odd characters.  The article's exclusions
of $\theta=1$ and $\theta=\omega^2$ are not used for the other
two characters.

Let $U_{\rm loc}$ be the inverse norm limit of local pro-$37$ units
and $E$ the image of the inverse limit of global units.  The actual
reciprocity maps give
\[
 0\longrightarrow E\longrightarrow U_{\rm loc}
 \longrightarrow M\longrightarrow X\longrightarrow0.
\]
The global odd units are exactly $E^-=\mathbb Z_{37}(1)$.
Indeed $v/c(v)$ is a root of unity for a global unit $v$, and
$2$ is invertible on its pro-$37$ completion.  The root-of-unity
inclusion and the first-moment map in Coleman's sequence are
\begin{equation}\label{eq:component-coleman}
 0\longrightarrow\mathbb Z_{37}(1)\longrightarrow U_{\rm loc}
 \xrightarrow{\operatorname{Col}}
 \mathbb Z_{37}[[\Delta\times\Gamma]]
 \xrightarrow{\lambda\mapsto\int x\,d\lambda(x)}
 \mathbb Z_{37}(1)\longrightarrow0.
\end{equation}
See \cite[Theorem~5.5.31]{SharifiIwa}; the prime-to-$37$
root term in the full sequence disappears on taking pro-$37$
units.  The class modules $X^{[-1]}$ and $X^{[1]}$ vanish, also
directly from \eqref{eq:component-p37-class-input}.  In character
$-1$, both ends of \eqref{eq:component-coleman} vanish and $E^{[-1]}=0$,
so reciprocity identifies $M^{[-1]}$ with $\Lambda$.  In character
$1$, reciprocity and \eqref{eq:component-coleman} identify
\[
 M^{[1]}\simeq U_{\rm loc}^{[1]}/\mathbb Z_{37}(1)
 \simeq\ker\bigl(\Lambda\longrightarrow\mathbb Z_{37}(1)\bigr)
 =(\gamma-u)\Lambda\simeq\Lambda,
 \qquad u=\kappa(\gamma).
\]
The kernel equality is division by the monic linear series
$T-(u-1)$, with $T=\gamma-1$.  Multiplication by $\gamma-u$ is
the displayed isomorphism from $\Lambda$ onto that ideal.  This
finishes the proof of \eqref{eq:component-p37-odd}.

The other input is a cup product at the first finite layer itself.
McCallum--Sharifi's calculation supplies $c$-even Kummer classes
\begin{equation}\label{eq:component-nonzero-cup}
 f,g\in H^1(P,\mathbb F_{37}(1)),\qquad
 f\smile g\ne0\quad\text{in }H^2(P,\mathbb F_{37}(2)).
\end{equation}
Here $P$ is the maximal pro-$37$ Galois group of
$F_1=\mathbb Q(\mu_{37})$ unramified outside $37$.
In their notation one can use the cyclotomic units $\eta_7$
and $\eta_{25}$: Corollary~8.3 has a nonzero coefficient of
$[x_7,x_{25}]$, and their duality formula identifies that
coefficient with this cup product
\cite[Theorem~7.5, p.~20; Corollary~8.3, p.~23]{MS}.
The unit characters are $\omega^{1-7}$ and $\omega^{1-25}$,
which are even.  The cup target has character $\omega^{2-32}$.
The degree-one classes for the full group $G_{F_1,S}$ factor
through its maximal pro-$37$ quotient $P$.  Naturality of cup
products under inflation shows that their nonzero cup for
$G_{F_1,S}$ comes from a nonzero cup for $P$; no surjectivity
assertion for degree-two inflation is required.
The two coefficient modules in \eqref{eq:component-nonzero-cup}
are trivial as $P$-modules because $\mu_{37}\subset F_1$.
Their $c$-actions are respectively $-1$ and $+1$.

\begin{theorem}\label{thm:component-no-galois-p37}
At $p=37$, let $t\in H$ be any lift of a
$\mathbb Z_{37}$-generator of $T=T_{32}$.  Its closed normal
closure $N=\langle t\rangle_H^{\rm normal}$ is not
$\Gamma$-stable.  The field $\Omega_{37}^N$ is therefore not
Galois over $F_1$, and hence is not Galois over $\mathbb Q$.
This applies in particular to every exact $\Delta$-eigenlift of
that generator.
\end{theorem}

\begin{proof}
The image of $t$ in $H/\Phi(H)=M/37M$ is nonzero, since it
generates the direct summand $T/37T$.  It is primitive in the free
pro-$37$ group $H$, so $H/N$ is free pro-$37$.
Here primitivity can be checked directly in countable rank.
Choose a convergent free basis $b_1,b_2,\ldots$ of $H$ and a
coordinate in which $t$ has nonzero Frattini coefficient, calling
it the first coordinate.  Replace $b_1$ by $t$ and leave the
others fixed.  The resulting endomorphism is an isomorphism on
the Frattini quotient, so its closed image is $H$: a proper
closed subgroup lies in an open subgroup of index $37$.
It preserves the normal subgroup $K_n$ generated by the tail
$b_j$, $j>n$, and its map on $H/K_n$ is a surjective
endomorphism of a finitely generated free pro-$37$ group, hence
an automorphism.  The groups $K_n$ have trivial intersection,
because any finite continuous quotient kills all but finitely
many elements of a convergent basis.  Thus the original
endomorphism is injective as well.  This verifies the free
factor containing $t$ and the asserted freeness of $H/N$.

The image of $N$ in $M$ is precisely $T$, whence
\begin{equation}\label{eq:component-quotient-ab}
 (H/N)^{\rm ab}=M/T=U\simeq\Lambda^{18}.
\end{equation}
Suppose for contradiction that $N$ is $\Gamma$-stable.  Then
$Q=P/N$ is defined, with the induced split extension
\[
 1\longrightarrow H/N\longrightarrow Q\longrightarrow\Gamma
 \longrightarrow1.
\]
Take $A=\mathbb F_{37}(2)$, which is a trivial $Q$-module.
Both $H/N$ and $\Gamma$ have $37$-cohomological dimension at
most one.  The Hochschild--Serre spectral sequence therefore
gives the actual isomorphism
\begin{equation}\label{eq:component-h2-zero}
 H^2(Q,A)\simeq
 H^1\bigl(\Gamma,\operatorname{Hom}_{\rm cts}(U,A)\bigr)=0.
\end{equation}
To verify the last equality including its topology, put
$V=\operatorname{Hom}_{\rm cts}(U,A)$, a discrete module.
Procyclic cohomology gives $H^1(\Gamma,V)=V/(\gamma-1)V$.
By \eqref{eq:component-quotient-ab}, the compact dual is
$U/37U\simeq\mathbb F_{37}[[T]]^{18}$.
The dual of $\gamma-1$ on $V$ is multiplication by
$\gamma^{-1}-1=-\gamma^{-1}T$ on this module.  It is injective,
so continuous Pontryagin duality makes $\gamma-1$ surjective
on $V$.  The Tate twist introduces no additional scalar here:
$\kappa(\gamma)^2=1$ in $\mathbb F_{37}^{\times}$.

Now take the two classes in \eqref{eq:component-nonzero-cup}.
Since $P$ acts trivially on $\mathbb F_{37}(1)$, their cocycles
are homomorphisms.  Being $c$-even means, with the coefficient
action included,
\[
 f(chc^{-1})=-f(h),\qquad g(chc^{-1})=-g(h)
 \quad(h\in P).
\]
The restrictions of $f$ and $g$ to $H$ factor through $M$.
The abelian image of $t$ lies in $T=M^+$, so $ctc^{-1}$ and $t$
have the same abelian image.  The displayed identities therefore
give $2f(t)=2g(t)=0$, and hence $f(t)=g(t)=0$.
Each homomorphism kills the whole $H$-normal
closure $N$.  They consequently descend to classes
$\bar f,\bar g\in H^1(Q,\mathbb F_{37}(1))$.
Naturality of cup products gives
\[
 f\smile g=\operatorname{inf}_{Q}^{P}(\bar f\smile\bar g)=0
\]
by \eqref{eq:component-h2-zero}, contradicting
\eqref{eq:component-nonzero-cup}.  Thus $N$ is not
$\Gamma$-stable.  The statements about fixed fields are now
exactly the normality criterion in Galois correspondence.
\end{proof}

Exact $c$-fixed lifts exist: for any lift $b$ of a positive abelian class,
$b(b^{-1}c(b))^{1/2}$ is $c$-fixed with the same abelian image.
The theorem proves a failure of invariance for every primitive
lift, not just the particular lift in
\eqref{eq:component-raw-lift}.  For a component prescription which
kills $T_{32}$, this is the entire nonzero even part at $37$.

\subsection{Ambient saturation and the precise remaining map}

Let $\mathcal G=\operatorname{Gal}(\Omega_p/\mathbb Q)$,
$\pi:H\to M$, and let $A\subset M$ be a closed
$\Delta\times\Gamma$-stable sum of components.  From a set of
lifts with abelian image generating $A$, form its $H$-normal
closure $N$ and then its ambient saturation
\[
 N^{\mathcal G}=
 \overline{\langle gng^{-1}:g\in\mathcal G,\ n\in N\rangle}.
\]
This is a closed $\mathcal G$-normal subgroup of $H$, and
\[
 \pi(N^{\mathcal G})=A,\qquad
 (H/N^{\mathcal G})^{\rm ab}=M/A.
\]
The equality of images follows by taking the closed sum of all
$\mathcal G$-translates of $A$.  Thus saturation repairs the
Galois action while preserving the specified abelian quotient.

\begin{corollary}\label{cor:component-saturation-not-free}
At $p=37$, let $S\subset H$ be any closed $P$-normal subgroup
whose image in $M$ is precisely $T$.  Then
\[
 H^2(P/S,\mathbb F_{37}(2))\ne0,
 \qquad H/S\text{ is not a free pro-}37\text{ group}.
\]
In particular this holds for the full ambient saturation of
the normal closure of any lift of a generator of $T_{32}$.
\end{corollary}

\begin{proof}
The two Kummer homomorphisms in
\eqref{eq:component-nonzero-cup} vanish on $S$, since their
restrictions to $H$ factor through odd functionals on $M$.
They descend to $P/S$, and their descended cup has nonzero
inflation.  It is therefore nonzero.  On the other hand,
$(H/S)^{\rm ab}=M/T=\Lambda^{18}$.
If $H/S$ were free pro-$37$, the Hochschild--Serre calculation
\eqref{eq:component-h2-zero}, with $S$ in place of $N$, would
make this same degree-two cohomology zero.  This contradiction
proves the second assertion.  Ambient saturation has image
$T$ by the preceding image calculation, so it is included.
\end{proof}

Put $K=\operatorname{Gal}(\Omega_p/\Omega_p^*)$ for the actual
outer-action kernel.  Galois correspondence gives the precise
remaining requirement
\[
 \Omega_p^{N^{\mathcal G}}\subseteq\Omega_p^*
 \quad\Longleftrightarrow\quad K\subseteq N^{\mathcal G}.
\]
The automatic intersection construction is
\[
 \Omega_p^{N^{\mathcal G}K}
   =\Omega_p^{N^{\mathcal G}}\cap\Omega_p^*,
 \qquad
 (H/(N^{\mathcal G}K))^{\rm ab}
   =M/(A+\pi(K)).
\]
These formulas describe the exact maps and retain the needed
nonabelian kernel containment as a separate assertion.
The next construction proves that injection on a retained free
factor supplies no such containment, even after full saturation.

\subsection{An operator-group counterexample after full saturation}

This subsection concerns the explicit groups constructed here.
The arithmetic result is Theorem~\ref{thm:component-no-galois-p37}.
Fix $p=37$, an element $\alpha\in1+37\mathbb Z_{37}$, a
cyclic group $\Delta=\langle\delta\rangle$ of order $36$, and
its faithful Teichm\"uller character $\omega$.

We use the following exact eigenprojection.  For a pro-$37$ group
with this $\Delta$-action, put $\beta=\omega(\delta)^e$ and
\[
 \mathcal E_e(v)=
 \left(\prod_{j=0}^{35}\delta^j(v)^{\beta^{-j}}\right)^{1/36},
 \qquad
 v^{(e)}=\lim_{n\to\infty}\mathcal E_e^n(v).
\]
The ordered product uses increasing $j$.
For completeness, if $\delta(v)v^{-\beta}$ belongs to the
$m$th closed lower central subgroup, its conjugates commute
modulo the $(m+1)$st subgroup.  In that quotient the displayed
averaging is the ordinary character idempotent.  It changes $v$
only in degree $m$ and makes the eigen-relation hold in degree
$m+1$.  Repeated averaging therefore converges in every finite
$37$-group quotient, since such a quotient is nilpotent.
It yields $\delta(v^{(e)})=(v^{(e)})^\beta$.
Products, unit powers and limits show that this construction is
functorial for continuous $\Delta$-equivariant homomorphisms.
This is also the eigenprojection of
\cite[Lemma~4]{Sharifi}.

Let $B$ be free pro-$37$ on $\gamma$ and $x_a$ for all odd
residues $a$ modulo $36$.  Map $B$ onto $\Gamma\simeq\mathbb Z_{37}$
by $\gamma\mapsto\gamma$ and $x_a\mapsto1$, and put
$J=\ker(B\to\Gamma)$.  Set
\[
 \delta(\gamma)=\gamma,\qquad
 \delta(x_a)=x_a^{\omega(\delta)^a}.
\]
These formulas define an automorphism of order $36$.
The group $J$ is free pro-$37$, and its abelianization has the
operator-module description
\[
 J^{\rm ab}\simeq\bigoplus_{a\ \mathrm{odd}}\Lambda(\omega^a).
\]
Here are the two maps proving this description.  The normal
generators $x_a$ give a surjection $\Lambda^{18}\to J^{\rm ab}$.
The homomorphism $B\to\Lambda^{18}\rtimes\Gamma$ sending
$x_a$ to the corresponding standard vector gives a map in the
opposite direction.  Their composite fixes all standard vectors,
so the initial surjection is an isomorphism.  This also proves
the stated $\Delta$-action.

Form the free pro-$37$ group $\mathcal F=J*\langle z\rangle$
and extend the commuting operators by
\[
 \delta(z)=z^{\omega(\delta)^{32}},\qquad
 \gamma^b(z)=z^{\alpha^b}\quad(b\in\mathbb Z_{37}).
\]
On $J$, the $\Gamma$-action is the conjugation action from $B$.
Thus
\[
 \mathcal F^{\rm ab}=U_0\oplus T_0,
 \qquad
 U_0=\bigoplus_{a\ \mathrm{odd}}\Lambda(\omega^a),\qquad
 T_0=\Lambda/(\gamma-\alpha)
\]
with $\Delta$-character $\omega^{32}$ on $T_0$.
Choose
\[
 w=[x_3,x_{29}]^{(32)}\in J,
 \qquad t=(zw)^{(32)}\in\mathcal F.
\]
Both $z$ and $t$ are exact $\Delta$-eigenlifts of the same
generator of $T_0$.  The retraction
\[
 \rho:\mathcal F\longrightarrow J,\qquad
 \rho|_J=\operatorname{id},\quad\rho(z)=1
\]
is $\Gamma\times\Delta$-equivariant.  It is injective on the
retained free factor $J$, and $\rho(t)=w$.
Let $\mathcal N$ be the closed normal closure of $t$ in
$\mathcal F\rtimes(\Gamma\times\Delta)$.

In $\operatorname{UT}_3(\mathbb F_{37})$ put
\[
 X=I+E_{12},\qquad Y=I+E_{23},\qquad
 Z=[X,Y]=I+E_{13}.
\]
Map $B$ to this group by
\[
 \gamma\mapsto1,\quad x_3\mapsto X,\quad x_{29}\mapsto Y,
 \quad x_a\mapsto1\quad(a\ne3,29),
\]
and restrict the map to $J$.  Extend it to
$q:\mathcal F\to\operatorname{UT}_3(\mathbb F_{37})$ by
$q(z)=Z^{-1}$.  The map is $\Gamma$-invariant: this is true
on $J$ by its construction, and on $z$ because
$\alpha^b\equiv1\pmod{37}$.  It is $\Delta$-equivariant for
the Heisenberg automorphism that scales $X,Y,Z$ with respective
characters $\omega^3,\omega^{29},\omega^{32}$.
Consequently $\ker q$ is stable under all the ambient operators.

The eigenprojection is already exact on the central element $Z$,
so functoriality gives
\[
 q(w)=Z,\qquad q(t)=(Z^{-1}Z)^{(32)}=1.
\]
Hence $\mathcal N\subseteq\ker q$.  The element $z$ supplies
the required certificate:
\begin{equation}\label{eq:component-saturation-kernel-failure}
 z\in\ker\rho,\qquad q(z)=Z^{-1}\ne1,
 \qquad z\notin\mathcal N.
\end{equation}
Thus $\ker\rho\nsubseteq\mathcal N$ after full ambient
saturation, with all displayed abelian modules, eigenlifts and
operator actions fixed.  The quotient
$(\mathcal F/\mathcal N)^{\rm ab}$ is nevertheless $U_0$,
by the saturation image calculation above.
Equation~\eqref{eq:component-saturation-kernel-failure} also
rules out a factorization of the quotient map through $\rho$:
its evaluation on $z$ would send $1$ to a nonidentity element.

% Suggested additional bibliography entries (do not duplicate a main-file entry):
% \bibitem{MS}
% W. G. McCallum and R. T. Sharifi,
% \emph{A cup product in the Galois cohomology of number fields},
% Duke Math. J. 120 (2003), 269--310.
% Author PDF: https://www.math.ucla.edu/~sharifi/publ/pairing.pdf
% Exact page locators above refer to this author PDF (dated December 5, 2003).
%
% \bibitem{SharifiCyclotomicNotes}
% R. T. Sharifi, \emph{Modular curves and cyclotomic fields},
% Theorem 2.3.18, printed p.19.
% https://www.math.ucla.edu/~sharifi/misc/aws-iwasawa.pdf
